\subsection{Cropland Protection Policy and the Requisition--Compensation Balance System}

To ensure food security and agricultural self-sufficiency, the Chinese government enforces a cropland protection policy that requires the total area of cropland to remain stable over time. This aggregate quantity constraint, commonly referred to as the ``cropland red line``, mandates that the national cropland area not decline, conditional on no deterioration in overall land quality (\cite{StateCouncil2006NationalLandUsePlan_EN}).

The implementation of the cropland protection policy relies on a quota-based system embedded in China’s administrative hierarchy. The central government first sets a national cropland retention target, which is subsequently allocated to provincial governments. Provincial authorities further disaggregate these targets to prefecture- and county-level jurisdictions. As a result, each local government is assigned a binding cropland quota that it is responsible for fulfilling.

Maintaining this aggregate cropland target becomes particularly challenging in the context of rapid industrialization and urban expansion, especially in jurisdictions with rapid urban expansion. To reconcile urban development with cropland protection, the central government introduced the \emph{requisition--compensation balance system}(\cite{StateCouncil1998LandAdminRegulations}). Under this policy, any reclamation of cropland for non-agricultural use must be offset by the replenishment of an equivalent amount of new cropland, thereby preserving cropland quantity at the aggregate level.

In practice, local governments are granted limited flexibility in meeting these quotas through spatial adjustment. When a county is unable to achieve the requisition--compensation balance within its own jurisdiction, it may adjust cropland quotas with neighboring counties within the same prefecture. If such an adjustment remains infeasible, cross-county compensation can be arranged within the province among areas with comparable resource endowments. Through these within-jurisdiction adjustments, the policy aims to maintain aggregate cropland quantity while accommodating heterogeneity in local land conversion pressures.

\begin{figure}[htbp]
\begin{center}
  \includegraphics[width=0.6\textwidth]{background_crop_1.png}
\caption{Changes in Cropland Share in China}
\label{fig:landuse}
\end{center}
\footnotesize Data Source: China Land Cover Data.
\end{figure}

\subsection{Policy Revision and Monitoring Tightening in 2017}

Cropland protection has been a long-standing policy priority in China since the early 2000s. The State Council has imposed binding targets on cropland retention, requiring each administrative jurisdiction to maintain cropland area above an assigned benchmark. In practice, however, enforcement of these targets remained weak prior to 2017, potentially due to monitoring difficulties and an underdeveloped accountability system. As illustrated in Figure~\ref{fig:landuse}, the national share of cropland continued to decline steadily between 2005 and 2017, despite the existence of formal protection targets.

In order to strengthen the implementation of the requisition--compensation balance system and reinforce the cropland red line, the State Council issued the \textit{Opinions on Strengthening Cropland Protection and Improving the Requisition--Compensation Balance System}\footnote{\cite{CPCStateCouncil2017farmland}} in 2017 (hereafter, the 2017 Opinions). The 2017 Opinions revised the existing cropland protection framework by enhancing responsibility assignment and tightening monitoring mechanisms. In particular, while principal leaders of local governments at all administrative levels had long been identified as the first responsible persons for cropland protection, the 2017 Opinions further strengthened this responsibility by explicitly linking cropland protection performance to official evaluations, including cadre performance assessments and assessments related to ecological civilization construction.

The revised policy was implemented through a top-down governance structure. As shown in Figure~\ref{fig:timeline}, following the issuance of the 2017 Opinions by the central government in January 2017, a formal assessment framework for provincial governments was promulgated in January 2018. Provincial governments subsequently adopted corresponding responsibility assessment measures for lower-level jurisdictions, thereby cascading the strengthened protection requirements downward through the administrative hierarchy.

The stated objective of the 2017 policy revision was to establish a new cropland protection mechanism that emphasizes both quantity and productive capacity. On the quantity dimension, national cropland area stabilized after 2017 and exhibited a modest recovery, as shown in Figure~\ref{fig:landuse}. On the quality and productive capacity dimension, however, outcomes appear less clear-cut. Although the central government introduced monitoring systems intended to track cropland quality, ex post evidence suggests limited effectiveness. Reports by official media in the post-2017 period document persistent concerns regarding newly created cropland, including low land quality, abandonment, and subsequent loss \citep{Wang2023balance}. These accounts often attribute such outcomes to misaligned incentives and insufficient coordination across cropland protection policies and implementation measures.

\begin{figure}
    \centering
    \includegraphics[width=0.85\linewidth]{timeline.png}
    \caption{Timeline}
    \label{fig:timeline}
\end{figure}

\subsection{Jiangsu Province}

This study focuses on Jiangsu Province rather than conducting a nationwide analysis. The objective of the paper is to examine land misallocation and inefficient land use arising from multitask pressures faced by local governments, specifically the tension between economic development and cropland protection. Although cropland protection targets are formally imposed nationwide, they do not constitute a binding constraint in all provinces. In regions with relatively slow urbanization, weak economic growth pressure, persistent population outflows, or abundant land reserves, local governments face limited trade-offs between urban expansion and cropland retention. In such settings, the multitask incentives central to this study are unlikely to arise in a meaningful way.

Jiangsu provides a particularly suitable empirical setting for studying these trade-offs. It is one of China’s most economically developed provinces, with among the highest levels of GDP per capita and a long history of rapid urbanization. Located in the Yangtze River Delta, Jiangsu is characterized by flat terrain and a high share of land suitable for cultivation, which has historically supported intensive agricultural production. At the same time, sustained industrialization and urban expansion generate strong and persistent demand for construction land. As a result, local officials in Jiangsu face especially acute conflicts between meeting economic development objectives and complying with binding cropland protection targets, a tension that is repeatedly emphasized in provincial and central government policy documents. 

%(\textcolor{blue}{TBD 江苏省土地利用总体规划})
